% -----------------------------------------------------------
\section{Righteousness, Righteous}
% -----------------------------------------------------------
	\label{RIGHTEOUSNESS}
	
% ----------------------
\subsection{See also}
% ----------------------

	\hyperref[LAW]{Law},

% ----------------------
\subsection{1828 American Dictionary of the English Language} % *1828
% ----------------------
	\label{RIGHTEOUSNESS-1828}
	
	% -----
	\subsubsection{Righteousness}
	% -----

		\begin{compactitem}
			\item[1] Purity of heart and rectitude of life; conformity of heart and life to the divine law. \textit{righteousness} as used in Scripture and theology, in which it is chiefly used, is nearly equivalent to holiness, comprehending holy principles and affections of heart, and conformity of life to the divine law. It includes all we call justice, honesty and virtue, with holy affections; in short, it is true religion.
			\item[2] Applied to God, the perfection or holiness of his nature; exact rectitude; faithfulness.
			\item[3] The active and passive obedience of Christ, by which the law of God is fulfilled.
			\item[4] Justice; equity between man and man.
			\item[5] The cause of our justification.
		\end{compactitem}

% % ----------------------
% \subsection{Oxford English Dictionary} % *OED
% % ----------------------
	% \label{RIGHTEOUSNESS-OED}

	% \begin{compactitem}
		% \item[]
	% \end{compactitem}

% ----------------------
\subsection{Scriptural Usage} % *SCRIPTURES
% ----------------------
	\label{RIGHTEOUSNESS-SCRIPTURES}

	% % -----
	% \subsubsection{Old Testament} % *OT
	% % -----
		% \label{RIGHTEOUSNESS-OT}
	
		% % --
		% \paragraph{KJV}
		% % --
			% \label{RIGHTEOUSNESS-OT-KJV}
	
			% \begin{tabular}{ll}
				% Total occurrences in the OT & 
			% \end{tabular}
			
			% \begin{compactdesc}
				% \item[]
			% \end{compactdesc}
			
		% % --
		% \paragraph{Hebrew Bible}
		% % --
			% \label{RIGHTEOUSNESS-HEBREW}
		
			% %
			% \subparagraph{\texorpdfstring{\he{sample word}}{}} % paste actual hebrew text here
			% %
				% \label{PAGA} % whatever is in step bible without periods
			
				% \begin{compactdesc}
					% \item[HALOT , PDF ] \textbf{} % Use the 1994 PDF
						% \begin{compactitem}
							% \item[1]
						% \end{compactitem}
					% \item[BDB , PDF ] \textbf{}
						% \begin{compactitem}
							% \item[1]
						% \end{compactitem}
					% \item[DCH PDF ] \textbf{}
						% \begin{compactitem}
							% \item[1]
						% \end{compactitem}
				% \end{compactdesc}
				
				% \begin{compactdesc}
					% \item[Some OT uses] \textbf{}
						% \begin{compactdesc}
							% \item[]
						% \end{compactdesc}
				% \end{compactdesc}
			
		% % --
		% \paragraph{Septuagint}
		% % --
			% \label{RIGHTEOUSNESS-LXX}
		
			% %
			% \subparagraph{\texorpdfstring{\gr{sample word}}{}}
			% %
		
	% -----
	\subsubsection{New Testament} % *NT
	% -----
		\label{RIGHTEOUSNESS-NT}
	
		% --
		\paragraph{KJV}
		% --
			\label{RIGHTEOUSNESS-NT-KJV}
			
	
			\begin{tabular}{ll}
				Total occurrences in the NT & 99
			\end{tabular}
			
			\begin{compactdesc}
				\item[{\hyperref[MATTHEW-5-20]{Matthew 5:20}}] \phantom{.}
					\begin{compactitem}
						\item Need to check out the Greek here and what alternate translations say.
					\end{compactitem}
				\item[{\hyperref[LUKE-15-4]{Luke 15:4--7}}]
				\item[{\hyperref[MATTHEW-6-33]{Matthew 6:33}}]
				\item[{\hyperref[ROMANS-3-20]{Romans 3:20--22}}]
				\item[{\hyperref[ROMANS-4]{Romans 4}}] \phantom{.}
					\begin{compactitem}
						\item An interesting little discourse on righteousness.
						\item Very important to work out the relationship between belief/faith and righteousness.
						\item I think this chapter is relevant to some issues in faith vs. works, like verses 4--5 for example.
					\end{compactitem}
				\item[{\hyperref[ROMANS-8-4]{Romans 8:4, 10}}] 
				\item[{\hyperref[2CORINTHIANS-6-14]{2 Corinthians 6:14}}] 
				\item[{\hyperref[GALATIANS-3-21]{Galatians 3:21}}] \phantom{.}
					\begin{compactitem}
						\item A very important verse for how Paul sees the relationship between righteousness and \textit{life}.
					\end{compactitem}
			\end{compactdesc}
			
		% --
		\paragraph{Greek New Testament} % *GREEK
		% --
			\label{RIGHTEOUSNESS-GREEK}
			
			%
			\subparagraph{\texorpdfstring{\gr{ἀδικία}}{}}
			%
				\label{ADIKIA}
			
				\begin{compactdesc}
					\item[BDAG 18] \textbf{}
						\begin{compactitem}
							\item[1] \textbf{an act that violates standards of right conduct, \textit{wrongdoing}} (opposite \gr{δικαιοσύνη})
							\item[2] \textbf{the quality of injustice, \textit{unrighteousness, wickedness, injustice}}...in a broader sense \gr{ἀδικία} can refer to the essence of the evil world...in such case the genitive signifies one's allegiance or subjugation to it
						\end{compactitem}
					\item[CGL 20a, PDF 46] \textbf{}
						\begin{compactitem}
							\item \textbf{injustice, wrongdoing}
						\end{compactitem}
					\item[LSJ 23b, PDF 96] \textbf{}
						\begin{compactitem}
							\item[I] \textit{wrongdoing, injustice}
							\item[II] \textit{wrongful act, offence}
						\end{compactitem}
				\end{compactdesc}
				
				\begin{tabular}{ll}
					Total occurrences in the NT & 25
				\end{tabular}
				
				% \begin{compactdesc}
					% \item[Some NT uses] \textbf{}
						% \begin{compactdesc}
							% \item[]
						% \end{compactdesc}
				% \end{compactdesc}
			
			%
			\subparagraph{\texorpdfstring{\gr{δίκαιος}}{}}
			%
				\label{DIKAIOS}
			
				\begin{compactdesc}
					\item[BDAG 218a] \textbf{}
						\begin{compactitem}
							\item[1] \textbf{pertaining to being in accordance with high standards of rectitude, \textit{upright, just, fair}}
							\item[1.a] of humans
							\item[1.a.$\alpha$] In Greco-Roman tradition a \gr{δίκαιος} person is one who upholds the customs and norms of behavior, including especially public service, that make for a well-ordered, civilized society
							\item[1.a.$\beta$] of things relating to human beings
							\item[1.b] of transcendent beings. Because of their privileged status as authority figures, the idea of fairness or equity is associated with such entities
							\item[1.b.$\alpha$] God
							\item[1.b.$\beta$] of Jesus who, as the ideal of an upright person is called simply \gr{δίκαιος}
							\item[2] the neuter denotes that which is \textbf{obligatory in view of certain requirements of justice, \textit{right, fair, equitable}}
						\end{compactitem}
					\item[CGL 378b, PDF 404] \textbf{}
						\begin{compactitem}
							\item[1] (of persons) observing the rules and customs of society, \textbf{civilised, well-mannered, decent}
							\item[2] (of perons, their attributes) having due regard to what is right and honourable, \textbf{righteous, honourable, upright, fair}
							\item[3] (of actions, speech, or similar) conforming with right or law, \textbf{right, lawful, just}
							\item[4] (of persons) \textbf{right, proper}, or \textbf{deserving} (with infinitive to do or experience something)
							\item[5] (generally, or an account) \textbf{strictly true, precise}
						\end{compactitem}
					\item[LSJ 429a, PDF 500] \textbf{}
						\begin{compactitem}
							\item[A.1] in Homer and all writers, of persons, \textit{observant of custom} or \textit{rule}
							\item[A.2] \textit{observant of duty} to gods and men, \textit{righteous}
							\item[A.3] euphemistic of a sacred snake
							\item[B.I.a] \textit{equal, even, well-balanced}
							\item[B.I.b.1] \textit{legally exact, precise}
							\item[B.I.b.2] \textit{lawful, just}
							\item[B.II.a] of persons and things, \textit{meet and right, fitting}
							\item[B.II.b.1] \textit{normal}
							\item[B.II.b.2] \textit{real, genuine} 
							\item[B.II.b.3] the plea of \textit{equity}
							\item[B.III] \textit{(special use)}
							\item[C] \textit{(special, apparently impersonal use in some prose sources)}
						\end{compactitem}
				\end{compactdesc}
				
				% \begin{tabular}{ll}
					% Total occurrences in the NT & 
				% \end{tabular}
				
				\begin{compactdesc}
					\item[Some NT uses] \textbf{}
						\begin{compactdesc}
							\item[{\hyperref[GALATIANS-3-11]{Galatians 3:11}}] \phantom{.}
								\begin{compactitem}
									\item There are other verses around here that talk about justification, which is the same root as \gr{δίκαιος}.
								\end{compactitem}
						\end{compactdesc}
				\end{compactdesc}
		
			%
			\subparagraph{\texorpdfstring{\gr{δικαιοσύνη}}{}}
			%
				\label{DIKAIOSUNE}
				
				See also \hyperref[ANOMIA]{\grl{ἀνομία}}; BDAG and LSJ give the opposite of \gr{ἀνομία} as \gr{δικαιοσύνη}, so this word is important in thinking about the opposite of lawlessness or lawless conduct (i.e., lawfulness and lawful conduct).
				
				BDAG also gives \gr{ἀδικία} as an opposite of \gr{δικαιοσύνη} (in the entry for the former).
			
				\begin{compactdesc}
					\item[BDAG 219a] generally, the quality of being upright. Theognis (sixth century BC) defines \gr{δικαιοσύνη} as the sum of all \gr{ἀρετή}; according to Demosthenes it is the opposite of \gr{κακία}. A strict classification of \gr{δικαιοσύνη} in the NT is complicated by frequent interplay of abstract and concrete aspects drawn from the OT and Greco-Roman cultures, in which a sense of equitableness combines with an awareness of responsibility within a social context.
						\begin{compactitem}
							\item[1] \textbf{the quality, state, or practice of judicial responsibility with focus on fairness, \textit{justice, equitableness, fairness}}
							\item[1a] of human beings
							\item[1b] of transcendent figures
							\item[2] \textbf{quality or state of juridical correctness with focus on redemptive action, \textit{righteousness}}. Equitableness if especially associated with God.
							\item[3] \textbf{the quality or characteristic of upright behavior, \textit{uprightness, righteousness}}
							\item[3a] of uprightness in general
							\item[3b] of specific action, \textit{righteousness} in the sense of fulfilling divine expectation not specifically expressed in ordinances
							\item[3c] uprightness as determined by divine legal standards
						\end{compactitem}
					\item[CGL 378b, PDF 404] \textbf{}
						\begin{compactitem}
							\item \textbf{righteousness, justice}
						\end{compactitem}
					\item[LSJ 429a, PDF 500] \textbf{}
						\begin{compactitem}
							\item[I.1] \textit{righteousness, justice}
							\item[I.2] \textit{fulfillment of the Law} \textit{(Here the entry mentions the LXX and the gospel of Matthew, and notes there are other uses in the NT as well.)}
							\item[II] \textit{justice, the business of a judge}
							\item[III] \gr{Δικαιοσύνη}, personified
							\item[IV] Pythagorean name for \textit{four}
							\item[V] \textit{(Some special use limited to a certain author.)}
						\end{compactitem}
				\end{compactdesc}
				
				% \begin{tabular}{ll}
					% Total occurrences in the NT & 
				% \end{tabular}
				
				\begin{compactdesc}
					\item[Some NT uses] \textbf{}
						\begin{compactdesc}
							\item[2 Corinthians 6:14] \gr{Μὴ γίνεσθε ἑτεροζυγοῦντες ἀπίστοις· τίς γὰρ μετοχὴ δικαιοσύνῃ καὶ ἀνομίᾳ, ἢ τίς κοινωνία φωτὶ πρὸς σκότος;}
							\item[{\hyperref[GALATIANS-3-21]{Galatians 3:21}}]
						\end{compactdesc}
				\end{compactdesc}
				
			%
			\subparagraph{\texorpdfstring{\gr{ἐλεημοσύνη}}{}}
			%
				\label{ELEEMOSUNE}
				
				% \begin{tabular}{ll}
					% Total occurrences in the NT & 
				% \end{tabular}
				
				% \begin{compactdesc}
					% \item[Some NT uses] \textbf{} % Change to "all" if they're all listed
						% \begin{compactdesc}
							% \item[]
						% \end{compactdesc}
				% \end{compactdesc}
			
				\begin{compactdesc}
				
					\item[BDAG 279b] \textbf{}
						\begin{compactitem}
							\item[1] \textbf{exercise of benevolent goodwill, \textit{alms, charitable giving}} with focus on attitude and action as such
							\item[2] \textbf{that which is benevolently given to meet a need, \textit{alms}} with focus on material as such
						\end{compactitem}
						
					% \item[Brill , PDF ] \textbf{}
						% \begin{compactitem}
							% \item 
						% \end{compactitem}
						
					% \item[CGL , PDF ] \textbf{}
						% \begin{compactitem}
							% \item 
						% \end{compactitem}
						
					% \item[LSJ , PDF ] \textbf{}
						% \begin{compactitem}
							% \item 
						% \end{compactitem}
						
				\end{compactdesc}
		
	% -----
	\subsubsection{Book of Mormon} % *BOM
	% -----
		\label{RIGHTEOUSNESS-BOM}
	
		% \begin{tabular}{ll}
			% Total occurrences in the BOM & 
		% \end{tabular}
		
		\begin{compactdesc}
			\item[{\hyperref[HELAMAN-7-20]{Helaman 7:20--24}}]
			\item[{\hyperref[HELAMAN-13-13]{Helaman 13:13--14}}]
			\item[{\hyperref[HELAMAN-13-38]{Helaman 13:38}}]
			\item[{\hyperref[ETHER-10-11]{Ether 10:11}}]
			\item[{\hyperref[ETHER-10-13]{Ether 10:13}}]
		\end{compactdesc}
		
	% -----
	\subsubsection{Doctrine and Covenants} % *DC
	% -----
		\label{RIGHTEOUSNESS-DC}
	
		% \begin{tabular}{ll}
			% Total occurrences in the DC & 
		% \end{tabular}
		
		\begin{compactdesc}
			\item[{\hyperref[DC67-9]{DC 67:9}}]
			\item[{\hyperref[DC90-24]{DC 90:24}}]
		\end{compactdesc}
		
	% % -----
	% \subsubsection{Pearl of Great Price} % *PGP
	% % -----
		% \label{RIGHTEOUSNESS-PGP}
	
		% \begin{tabular}{ll}
			% Total occurrences in the PGP & 
		% \end{tabular}
		
		% \begin{compactdesc}
			% \item[]
		% \end{compactdesc}
		
% ----------------------
\subsection{Key Phrases} % *KEYPHRASES *PHRASES
% ----------------------
	\label{RIGHTEOUSNESS-PHRASES}

	\begin{compactdesc}
	
		\item[robe of righteousness] \textbf{5} | \hyperref[ISAIAH-61-10]{Isaiah 61:10}; \hyperref[2NEPHI-4-33]{2 Nephi 4:33}; \hyperref[2NEPHI-9-14]{2 Nephi 9:14}; \hyperref[DC29-12]{DC 29:12}; \hyperref[DC109-76]{DC 109:76}
			\begin{compactdesc}
				\item[Bible anchor verses] \phantom{.}
					\begin{compactdesc}
						\item[Isaiah 61:10] \he{m:'iyl .s:dAqAh}
					\end{compactdesc}
				\item[Commentary] Note also \hyperref[JOB-29-14]{Job 29:14}.
			\end{compactdesc}
			
	\end{compactdesc}
	
% % ----------------------
% \subsection{Bible Dictionaries} % *DICTIONARIES
% % ----------------------
	% \label{RIGHTEOUSNESS-BD}

% ----------------------
\subsection{Restoration Usage} % *RESTORATION
% ----------------------
	\label{RIGHTEOUSNESS-RESTORATION}

	% -----
	\subsubsection{Joseph Smith} % *JS
	% -----
		\label{RIGHTEOUSNESS-JS}
	
		\begin{compactdesc}
			\item[{\hyperlink{EMS-1843-JS-21-MAY-1843-HC-c}{21 May 1843 HC-c}}] \hl{Righteousness is not that which men esteem holiness. That which the world call righteousness I have not any regard for. To be righteous is to be just and merciful. if a man fails in kindness justice and mercy he will be damed} for many will say in that day Lord. have we not Prophesied in thy name and in thy name done many wonderful works but he will say unto them ye workers of iniquity \&c
		\end{compactdesc}
		
	% % -----
	% \subsubsection{Temple Ordinances}
	% % -----
		% \label{RIGHTEOUSNESS-TEMPLE}
	
		% \begin{compactdesc}
			% \item[{\hyperref[KEY]{Date/Source}}]
		% \end{compactdesc}
	
	% % -----
	% \subsubsection{Brigham Young} % *BY
	% % -----
		% \label{RIGHTEOUSNESS-BY}
	
		% \begin{compactdesc}
			% \item[{\hyperref[KEY]{Date/Source}}]
		% \end{compactdesc}
	
% % ----------------------
% \subsection{Commentary and Studies} % *COMMENTARY *STUDIES
% % ----------------------
	% \label{RIGHTEOUSNESS-COMMENTARY}

	% % -----
	% \subsubsection{Key Points}
	% % -----
		% \label{RIGHTEOUSNESS-KEYS}
	
		% \begin{compactitem}
			% \item
		% \end{compactitem}
		
	% % -----
	% \subsubsection{Sample Study}
	% % -----
		% \label{RIGHTEOUSNESS-study}